\subsection{Excepciones Encadenadas}

\subsubsection{Regla: preservar la excepción original como causa}

\textbf{Regla:} Cuando una excepción se traduce a otro tipo porque la capa actual puede aportar contexto útil, se conserva la excepción original como \texttt{innerException}. Se prefiere un tipo predefinido de .NET y solo se utiliza una excepción personalizada cuando el llamador necesita distinguir una condición propia del dominio (Microsoft, 2025a).

\textbf{Código correcto:}
\begin{lstlisting}[style=csharpstyle]
try
{
    _board.SaveState(currentState);
}
catch (IOException ex)
{
    throw new InvalidOperationException("Failed to persist game state.", ex);
}
\end{lstlisting}

\textbf{Código incorrecto:}
\begin{lstlisting}[style=csharpstyle]
try
{
    _board.SaveState(currentState);
}
catch (IOException ex)
{
    throw new InvalidOperationException("Failed to persist game state: " + ex.Message);
}
\end{lstlisting}

\subsubsection{Regla: \texttt{throw;} en lugar de \texttt{throw ex;}}

\textbf{Regla:} Al relanzar la misma excepción capturada, se usa \texttt{throw;} sin especificar la variable (Microsoft, 2026a).

\textbf{Descripción:} \texttt{throw ex;} reinicia la traza de la pila; \texttt{throw;} conserva la original completa (Microsoft, 2025a).

\textbf{Código correcto:}
\begin{lstlisting}[style=csharpstyle]
try
{
    _board.MovePlayer(player, destination);
}
catch (InvalidMoveException ex)
{
    _logger.Error($"Move failed. Player={player.Name}.", ex);
    throw;
}
\end{lstlisting}

\textbf{Código incorrecto:}
\begin{lstlisting}[style=csharpstyle]
try
{
    _board.MovePlayer(player, destination);
}
catch (InvalidMoveException ex)
{
    _logger.Error($"Move failed. Player={player.Name}.", ex);
    throw ex;
}
\end{lstlisting}
